\documentclass[aspectratio=169,11pt]{beamer}
\usepackage[T1]{fontenc}
\usepackage[utf8]{inputenc}
\usepackage{lmodern,amsmath,amssymb,mathtools,bm,booktabs,multicol,graphicx,tikz,xcolor,colortbl}
\usepackage{microtype}
\pdfstringdefDisableCommands{\let\translate\relax}
\usetikzlibrary{arrows.meta,positioning,calc,shapes.geometric}
\definecolor{paper}{HTML}{FBFAF7}
\definecolor{ink}{HTML}{282A2B}
\definecolor{muted}{HTML}{666965}
\definecolor{accent}{HTML}{945E49}
\definecolor{rulecolor}{HTML}{DCD9D1}
\definecolor{soft}{HTML}{F1EFE9}
\definecolor{pale}{HTML}{EEE8DF}
\setbeamercolor{background canvas}{bg=paper}
\setbeamercolor{normal text}{fg=ink,bg=paper}
\setbeamercolor{frametitle}{fg=ink,bg=paper}
\setbeamercolor{title}{fg=ink}
\setbeamercolor{structure}{fg=accent}
\setbeamercolor{alerted text}{fg=accent}
\setbeamercolor{block title}{fg=accent,bg=paper}
\setbeamercolor{block body}{fg=ink,bg=paper}
\setbeamercolor{block title example}{fg=ink,bg=paper}
\setbeamercolor{block body example}{fg=ink,bg=paper}
\setbeamerfont{frametitle}{size=\large,series=\bfseries}
\setbeamerfont{title}{size=\LARGE,series=\bfseries}
\setbeamerfont{block title}{series=\bfseries,size=\normalsize}
\setbeamertemplate{navigation symbols}{}
\setbeamertemplate{headline}{}
\setbeamertemplate{blocks}[default]
\setbeamertemplate{itemize item}{\raisebox{.22ex}{\tikz\fill[accent] (0,0) circle(1.25pt);}}
\setbeamertemplate{itemize subitem}{\raisebox{.22ex}{\tikz\fill[muted] (0,0) circle(1pt);}}
\setbeamertemplate{enumerate item}{\color{accent}\insertenumlabel.}
% Titles get room to breathe; remove the saturated full-width underline.
\setbeamertemplate{frametitle}{%
  \vspace{.18cm}\hspace{.41cm}\insertframetitle\par\vspace{.15cm}%
}
% Minimal footer: no repetitive journal citation on every slide.
\setbeamertemplate{footline}{%
\leavevmode\hbox{%
  \begin{beamercolorbox}[wd=.82\paperwidth,ht=2.7ex,dp=1.1ex,leftskip=.48cm]{normal text}%
    {\fontsize{7.5}{9}\selectfont\color{muted}SUPERFLEX\quad /\quad GEOMETRIC DEEP LEARNING}%
  \end{beamercolorbox}%
  \begin{beamercolorbox}[wd=.18\paperwidth,ht=2.7ex,dp=1.1ex,rightskip=.48cm plus1fil]{normal text}%
    {\fontsize{8}{9}\selectfont\color{muted}\insertframenumber}%
  \end{beamercolorbox}%
}}
% Section indicator in the upper-right corner
\setbeamertemplate{headline}{%
  \ifnum\value{section}>0
    \leavevmode
    \hbox to \paperwidth{%
      \hfill
      {\fontsize{8}{9}\selectfont\bfseries\color{accent}%
       \thesection\hspace{.12cm}\insertsectionhead}%
      \hspace{.45cm}%
    }%
    \vspace{-.05cm}%
  \fi
}
\setbeamersize{text margin left=0.43cm,text margin right=0.43cm}
\arrayrulecolor{rulecolor}
\newcommand{\caccent}[1]{{\color{accent}\bfseries #1}}
\newcommand{\ctiny}[1]{{\fontsize{8}{9}\selectfont\color{muted}#1}}
\newcommand{\src}[1]{\vspace{0.015cm}{\fontsize{7.1}{8.5}\selectfont\color{muted}#1}}
\newcommand{\eqbox}[1]{\begin{beamercolorbox}[sep=0.13cm]{block body}\centering\displaystyle #1\end{beamercolorbox}}
\newcommand{\safeincludegraphics}[2][]{%
\IfFileExists{#2}{\includegraphics[#1]{#2}}{%
\fcolorbox{rulecolor}{soft}{\parbox[c][2.0cm][c]{0.72\linewidth}{\centering\footnotesize\color{muted}Figure unavailable in this copy}}%
}}
\newcommand{\img}[2]{\safeincludegraphics[width=#1]{assets/#2}}
\newcommand{\timemark}[1]{\hfill{\color{muted}\footnotesize #1}}
\title{SuperFlex}
\subtitle{Deformable superquadrics for point-cloud decomposition}
\author{Geometric Deep Learning | Paper presentation}
\date{20-minute talk + 10-minute discussion}


\begin{document}


% 1
\begin{frame}[plain]
\vspace{0.40cm}
{\color{accent}\fontsize{10.5}{12}\selectfont\bfseries GEOMETRIC DEEP LEARNING}\par\vspace{0.28cm}
{\color{ink}\fontsize{39}{42}\selectfont\bfseries SuperFlex}\par\vspace{0.10cm}
{\color{muted}\large Deformable Superquadrics for Point Cloud Decomposition}\par\vspace{0.25cm}
{\fontsize{9}{11}\selectfont Gabriel Tavernini, Elisabetta Fedele, Tiago Novello, Leonidas Guibas, Marc Pollefeys, Francis Engelmann}\par
{\fontsize{9}{11}\selectfont\color{muted}ECCV 2026}\par\vspace{0.24cm}
\begin{center}\safeincludegraphics[width=.79\linewidth,height=2.45cm,keepaspectratio]{assets/paper_teaser.png}\end{center}
\src{Tavernini et al., ECCV 2026, Fig. 1.}
\end{frame}
% Contents
\begin{frame}{Contents}
\vspace{.35cm}
\begingroup
\setbeamerfont{section in toc}{size=\normalsize,series=\bfseries}
\setbeamertemplate{section in toc}{\color{accent}\inserttocsectionnumber\hspace{.2cm}\color{ink}\inserttocsection\par\vspace{.25cm}}
\tableofcontents[hideallsubsections]
\endgroup
\end{frame}


% 2
\section{Motivation and previous work}
\subsection{The decomposition problem}
\begin{frame}[shrink=9]{Understanding a 3D object through a few geometric parts}
\begin{columns}[T,onlytextwidth]
\column{.52\textwidth}
\vspace{.2cm}
\begin{block}{Input}
The input is a point cloud $\mathcal P=\{\mathbf x_i\}_{i=1}^N$ with $\mathbf x_i\in\mathbb R^3$.
\end{block}
\begin{block}{Output}
The output is a compact set of parametric solids whose union approximates the object through meaningful geometric parts.
\end{block}
\begin{itemize}\setlength\itemsep{.20em}
\item \textbf{Fidelity} requires a close match to the target geometry.
\item \textbf{Parsimony} favours a small number of primitives.
\item \textbf{Robustness} matters when parts of the surface are missing.
\end{itemize}
\column{.46\textwidth}
\centering\img{\linewidth}{paper_qualitative_shapenet.png}
\src{Paper Fig. 4. Each colour identifies a distinct primitive.}
\end{columns}
\end{frame}


% 3
\subsection{Previous approaches}
\begin{frame}{From earlier methods to SuperFlex}
\vspace{.15cm}
\begin{columns}[T,onlytextwidth]
\column{.33\textwidth}
\begin{block}{Optimization-based}
\textbf{EMS} (2022) introduced probabilistic fitting with improved robustness to outliers.\\[.14cm]
\textbf{Marching Primitives} (2023) uses iterative fitting guided by signed distance fields.
\end{block}
\column{.33\textwidth}
\begin{block}{Learning-based}
\textbf{SQ} (2019) learned superquadric shape abstractions.\\[.14cm]
\textbf{CSA} (2021) jointly estimated cuboids and point assignments.\\[.14cm]
\textbf{SuperDec} (2025) introduced Transformer-based set prediction.
\end{block}
\column{.33\textwidth}
\begin{block}{SuperFlex (2026)}
Inherits SuperDec's point-encoder + Transformer decoder.\\[.14cm]
Adds expressive geometry, stronger losses, optional refinement and occlusion training.
\end{block}
\end{columns}
\end{frame}


% 4
\section{Superquadric geometry}
\subsection{Canonical primitives}
\begin{frame}[shrink=6]{Superellipsoids}
\begin{columns}[T,onlytextwidth]
\column{.53\textwidth}
\vspace{.22cm}
{\large Canonical implicit surface}\par\vspace{.13cm}
\begin{equation*}\boxed{\displaystyle
q_{\mathrm{can}}(\mathbf x)=
\left[\left|\frac{x}{s_x}\right|^{\frac{2}{\epsilon_2}}+
\left|\frac{y}{s_y}\right|^{\frac{2}{\epsilon_2}}\right]^{\frac{\epsilon_2}{\epsilon_1}}
+\left|\frac{z}{s_z}\right|^{\frac{2}{\epsilon_1}} }
\end{equation*}
\vspace{.0cm}\centering\Large $\mathcal S=\{\mathbf x:q_{\rm can}(\mathbf x)=1\}$
\vspace{.16cm}\par\raggedright
\begin{itemize}\setlength\itemsep{.42em}
\item $\mathbf s=(s_x,s_y,s_z)$ sets the three axis scales.
\item $\epsilon_1$ controls the meridional profile.
\item $\epsilon_2$ controls the horizontal cross-section.
\end{itemize}
\column{.45\textwidth}
\centering\safeincludegraphics[width=\linewidth,height=4.5cm,keepaspectratio]{assets/demo_superquadrics.png}
\end{columns}
\end{frame}


% 5
\begin{frame}{Five canonical shape parameters are not enough}
\begin{columns}[T,onlytextwidth]
\column{.58\textwidth}
\centering\safeincludegraphics[width=\linewidth,height=4.9cm,keepaspectratio]{assets/demo_composition.png}
\column{.40\textwidth}
\begin{block}{Failure of rigid primitives}
Straight, symmetric parts are easy to model. Curved or tapered parts, however, can require several rigid primitives even when they correspond to a single meaningful component.
\end{block}
\vspace{.12cm}\begin{itemize}\setlength\itemsep{.5em}
\item \textbf{Tapering} changes the cross-sectional width along an axis.
\item \textbf{Bending} produces a curved central axis.
\end{itemize}
\end{columns}
\end{frame}


% 6
\subsection{Tapering and bending}
\begin{frame}{Tapering and bending}
\begin{columns}[T,onlytextwidth]
\column{.5\textwidth}
\centering\textbf{Tapering along $z$}\par\vspace{.05cm}
\safeincludegraphics[width=\linewidth,height=2.7cm,keepaspectratio]{assets/demo_taper.png}
\vspace{.01cm}\begin{equation*}\begin{aligned}
x'&=x(1+\tau_x z/s_z),\\[-1pt]
y'&=y(1+\tau_y z/s_z),\quad z'=z.
\end{aligned}\end{equation*}
\column{.49\textwidth}
\centering\textbf{Bending about $z$}\par\vspace{.05cm}
\safeincludegraphics[width=\linewidth,height=2.7cm,keepaspectratio]{assets/demo_bend.png}
\vspace{.01cm}\begin{itemize}\setlength\itemsep{.22em}
\item $k_z$ controls curvature, with circular radius $1/|k_z|$.
\item $\alpha_z$ sets the orientation of the bending plane.
\item Repeat for bending about $x$ and $y$.
\end{itemize}
\end{columns}
\end{frame}


% 7
\begin{frame}{Nineteen parameters, one implicit field}
\vspace{.09cm}
\begin{center}
\begin{tabular}{lccccc}
\toprule
Pose & Scale & Shape & Taper & Bend & \textbf{Total}\\
$(\mathbf t,R)$ & $\mathbf s$ & $(\epsilon_1,\epsilon_2)$ & $(\tau_x,\tau_y)$ & $\bm\beta=(k_a,\alpha_a)_{a=x,y,z}$ & \\
$3+3$ & $3$ & $2$ & $2$ & $6$ & \textbf{19} \\
\bottomrule
\end{tabular}
\end{center}
\vspace{.24cm}
\begin{equation*}
\boxed{q(\mathbf x;\Theta)=q_{\rm can}\!\left(
\underbrace{T^{-1}\circ B_y^{-1}\circ B_x^{-1}\circ B_z^{-1}}_{\text{undo deformations}}\bigl(R^\top(\mathbf x-\mathbf t)\bigr);\mathbf s,\bm\epsilon\right)}
\end{equation*}
\begin{itemize}\setlength\itemsep{.32em}
\item Evaluating a world-space point requires undoing its pose and geometric deformations.
\item On the surface $q(\mathbf x;\Theta)=1$; inside $q<1$ (in the valid invertible regime).
\end{itemize}
\end{frame}


% 8
\subsection{Radial distance}
\begin{frame}[shrink=3]{Radial distance as an efficient geometric proxy}
\begin{columns}[T,onlytextwidth]
\column{.54\textwidth}
\centering\safeincludegraphics[width=\linewidth,height=4.65cm,keepaspectratio]{assets/demo_radial.png}
\column{.44\textwidth}
\begin{block}{For the rigid case}
$\displaystyle q(\lambda\mathbf x)=\lambda^{2/\epsilon_1}q(\mathbf x)$\\[.2cm]
The surface is reached when $q(\lambda_{\ast}\mathbf x)=1$, giving\\[-.03cm]
$\displaystyle \lambda_{\ast}=q(\mathbf x)^{-\epsilon_1/2}$
\end{block}
\boxed{\displaystyle d_{rad}=\|R^\top(\mathbf x-\mathbf t)\|\bigl(1-q(\mathbf x;\Theta)^{-\epsilon_1/2}\bigr)}
\begin{itemize}\setlength\itemsep{.18em}
\item The sign distinguishes exterior points ($d>0$) from interior points ($d<0$).
\item This efficient proxy is not an exact Euclidean SDF.
\end{itemize}
\end{columns}
\end{frame}


% 9
\section{Architecture and learning}
\subsection{Primitive prediction}
\begin{frame}[shrink=8]{The model architecture}
\begin{center}\safeincludegraphics[width=.96\linewidth,height=3.06cm,keepaspectratio]{assets/paper_architecture.png}\end{center}
\vspace{.07cm}
\begin{columns}[T,onlytextwidth]
\column{.51\textwidth}
\begin{itemize}\setlength\itemsep{.28em}
\item The input point cloud is $\mathcal P\in\mathbb R^{N\times3}$.
\item A PVCNN encoder extracts point features $F_{PC}\in\mathbb R^{N\times D}$.
\item $P$ positional queries represent candidate superquadrics, with $F_{SQ}\in\mathbb R^{P\times D}$.
\end{itemize}
\column{.47\textwidth}
\begin{itemize}\setlength\itemsep{.28em}
\item $L$ Transformer layers combine query self-attention and point-feature cross-attention.
\item Prediction heads return $\{(\Theta_p,e_p)\}_{p=1}^P$ and soft point assignments.
\end{itemize}
\end{columns}
\vspace{.06cm}\centering
$\displaystyle W=\operatorname{softmax}_{p}\!\bigl(F_{PC}\,\phi(F_{SQ})^\top\bigr)\in\mathbb R^{N\times P}$
\end{frame}


% 10
\subsection{Reconstruction objective}
\begin{frame}{Why Chamfer distance alone is insufficient}
\begin{columns}[T,onlytextwidth]
\column{.53\textwidth}
\centering\safeincludegraphics[width=\linewidth,height=4.5cm,keepaspectratio]{assets/paper_qualitative_shapenet.png}
\src{Paper Fig. 4. SuperDec compared with rigid and deformable SuperFlex.}
\column{.45\textwidth}
\begin{block}{Chamfer distance}
Chamfer distance aligns samples on the observed and predicted surfaces, without directly comparing the volumes of their unions.
\end{block}
\begin{itemize}\setlength\itemsep{.52em}
\item Shape boundaries can become overly rounded.
\item Adjacent primitives can leave gaps.
\item Global volume and local surface accuracy are not jointly constrained.
\end{itemize}
\vspace{.13cm}\centering\caccent{The proposed loss combines global occupancy and local surface alignment.}
\end{columns}
\end{frame}


% 11
\begin{frame}[shrink=7]{Global volumetric consistency}
\begin{columns}[T,onlytextwidth]
\column{.55\textwidth}
\vspace{.13cm}
\begin{enumerate}\setlength\itemsep{.57em}
\item Convert the radial-distance proxy into \textbf{soft occupancy}\\[.05cm]
$\displaystyle I_p(\mathbf x)=\sigma\!\left(-\frac{d_p(\mathbf x)}{T_{\rm IoU}}\right)$.
\item Combine primitive occupancies with a \textbf{probabilistic OR}\\[.05cm]
$\displaystyle\hat O(\mathbf x)=1-\prod_{p=1}^P(1-e_p I_p(\mathbf x))$.
\item Compare with ground-truth occupancy $O$ over uniform samples $\mathcal V$.
\end{enumerate}
\column{.43\textwidth}
\centering\safeincludegraphics[width=\linewidth,height=3.25cm,keepaspectratio]{assets/demo_occupancy.png}
\vspace{.06cm}
{\footnotesize For $e_1=e_2=1$ we obtain\par}
$\hat O=I_1+I_2-I_1 I_2$.
\end{columns}
\vspace{.11cm}
\begin{equation*}
\boxed{\mathcal L_{\rm IoU}=-\log\frac{\sum_{\mathbf x\in\mathcal V}\hat O(\mathbf x)O(\mathbf x)}{\sum_{\mathbf x\in\mathcal V}[\hat O(\mathbf x)+O(\mathbf x)-\hat O(\mathbf x)O(\mathbf x)]}}
\end{equation*}
\end{frame}


% 12
\begin{frame}{Local surface alignment}
\begin{columns}[T,onlytextwidth]
\column{.59\textwidth}
\vspace{.18cm}
\begin{equation*}
\boxed{\mathcal L_{\rm SDF}=\frac1N\sum_{i=1}^N\sum_{p=1}^P w_{ip}\,\psi\bigl(d(\mathbf x_i;\Theta_p)\bigr)}
\end{equation*}
\begin{itemize}\setlength\itemsep{.47em}
\item $w_{ip}$ gives the soft assignment of point $i$ to primitive $p$.
\item $d$ is a signed radial \textbf{surrogate}, not the Euclidean SDF.
\item $\psi$ provides a bounded, asymmetric penalty that limits outlier influence.
\end{itemize}
\vspace{.12cm}
\column{.39\textwidth}
\centering\safeincludegraphics[width=\linewidth,height=3.85cm,keepaspectratio]{assets/demo_psi.png}
\vspace{.03cm}
{\footnotesize $\psi(d)=\begin{cases}\delta\tanh(d/\delta),& d\ge0,\\-\omega\delta\tanh(d/\delta),& d<0.\end{cases}$}
\end{columns}
\end{frame}


% 13
\begin{frame}{Compact decompositions require explicit regularization}
\vspace{.06cm}
\begin{equation*}
\boxed{\mathcal L=\lambda_{\rm IoU}\mathcal L_{\rm IoU}+\lambda_{\rm SDF}\mathcal L_{\rm SDF}
+\lambda_o\mathcal L_{\rm overlap}+\lambda_s\mathcal L_{\rm sparsity}+\lambda_e\mathcal L_{\rm exist}}
\end{equation*}
\vspace{.14cm}
\begin{columns}[T,onlytextwidth]
\column{.32\textwidth}
\begin{block}{Overlap}
$\displaystyle\mathbb E_{\mathbf x}\!\left[\operatorname{ReLU}(\sum_p I_p(\mathbf x)-1)\right]$\\[.1cm]
Penalizes excessive overlap between primitives.
\end{block}
\column{.32\textwidth}
\begin{block}{Sparsity}
$\displaystyle m_p=\frac1N\sum_i w_{ip}$\\[.1cm]
Discourages unnecessary primitives and favours compact decompositions.
\end{block}
\column{.32\textwidth}
\begin{block}{Existence}
$\displaystyle\hat e_p=\mathbf1[m_p>\epsilon_e]$\\[.1cm]
Binary cross-entropy aligns each existence prediction with primitive usage.
\end{block}
\end{columns}
\end{frame}


% 14
\section{Refinement and partial observations}
\subsection{Test-time refinement}
\begin{frame}{Optional test-time optimization}
\begin{columns}[T,onlytextwidth]
\column{.50\textwidth}
\begin{enumerate}\setlength\itemsep{.56em}
\item Select the primitives whose predicted existence satisfies $e_p>0.5$.
\item Optimize their geometry $\Theta_p$ by Adam, holding the set fixed.
\item Replace pointwise assignments with a global distance to the union.
\end{enumerate}
\vspace{.15cm}
\begin{equation*}\boxed{\widetilde d(\mathbf x)=-T\log\!\sum_{p\in\mathcal A}\exp\!\left(-\frac{d_p(\mathbf x)}{T}\right)}\end{equation*}
\centering{\footnotesize Smooth approximation to $\min_{p\in\mathcal A}d_p(\mathbf x)$.}
\column{.47\textwidth}
\centering\safeincludegraphics[width=\linewidth,height=3.75cm,keepaspectratio]{assets/paper_optimization.png}
\src{Paper Fig. 5. Refines existing primitives rather than training a new network.}
\end{columns}
\vspace{.15cm}\centering\caccent{IoU rises from 0.72 to 0.87, at a cost of 0.0082 s versus 5 s per object.}
\end{frame}


% 15
\subsection{Occlusion-aware training}
\begin{frame}{Partial scans}
\begin{columns}[T,onlytextwidth]
\column{.48\textwidth}
\centering\safeincludegraphics[width=\linewidth,height=4.4cm,keepaspectratio]{assets/paper_partial_ase.png}
\src{Paper Fig. 6. Partial inputs and reconstructions from standard and robust SuperFlex.}
\column{.50\textwidth}
\vspace{.08cm}
\begin{enumerate}\setlength\itemsep{.53em}
\item On complete clouds, predict and optimize high-quality primitive sets.
\item Treat the resulting sets as pseudo-ground truth.
\item Fine-tune on occluded inputs to reconstruct complete primitive structure.
\end{enumerate}
\vspace{.12cm}
\begin{beamercolorbox}[rounded=true,sep=.15cm]{block body}
Hungarian matching resolves the arbitrary ordering of the predicted primitives before Chamfer-based supervision is applied to corresponding surfaces.
\end{beamercolorbox}
\end{columns}
\end{frame}


% 16
\section{Experimental evaluation}
\subsection{Evaluation setup}
\begin{frame}{Experimental protocol}
\vspace{.11cm}
\begin{columns}[T,onlytextwidth]
\column{.49\textwidth}
\begin{block}{Object-level evaluation on ShapeNet}
\begin{itemize}\setlength\itemsep{.35em}
\item 13 object categories.
\item 4096 sampled points/object.
\item 16 maximum primitive queries.
\item 1000 training epochs.
\item 4 $\times$ NVIDIA GH200 GPUs reported for training.
\end{itemize}
\end{block}
\column{.49\textwidth}
\begin{block}{Partial observations}
\begin{itemize}\setlength\itemsep{.35em}
\item ASE provides partial observations with complete geometry for quantitative evaluation.
\item ScanNet++ supplies real-world scenes for qualitative evaluation.
\item Metrics include IoU, F-score, Chamfer L1/L2, primitive count and runtime.
\end{itemize}
\end{block}
\end{columns}
\vspace{.25cm}\centering\caccent{Separate feed-forward inference from the much slower optional refinement.}
\end{frame}


% 17
\subsection{Accuracy and ablations}
\begin{frame}{Quantitative comparison for ShapeNet}
\begin{center}
\renewcommand{\arraystretch}{1.25}
\begin{tabular}{lrrrr}
\toprule
Method & IoU $\uparrow$ & F $\uparrow$ & \# primitives $\downarrow$ & Runtime $\downarrow$ \\
\midrule
Marching Primitives & 0.56 & 0.18 & 24.10 & 163.41 s\\
SuperDec & 0.59 & 0.28 & 5.79 & 0.0075 s\\
\textbf{SuperFlex (rigid)} & \textbf{0.70} & 0.35 & 6.04 & 0.0075 s\\
\rowcolor{pale}\textbf{SuperFlex (deformable)} & \textbf{0.72} & \textbf{0.37} & \textbf{5.64} & 0.0082 s\\
\bottomrule
\end{tabular}
\end{center}
\vspace{.2cm}
\begin{columns}[T,onlytextwidth]
\column{.49\textwidth}
\begin{block}{Largest advance}
With rigid primitives, the new loss raises IoU from 0.59 to 0.70.
\end{block}
\column{.49\textwidth}
\begin{block}{Added expressiveness}
Tapering and bending raise IoU from 0.70 to 0.72 while preserving a compact representation.
\end{block}
\end{columns}
\src{Paper Table 1, ShapeNet. Forward passes only; all numbers are paper-reported, not reproduced here.}
\end{frame}


% 18
\begin{frame}{Which loss terms matter?}
\vspace{.04cm}
\begin{center}\renewcommand{\arraystretch}{1.22}
\begin{tabular}{lrrrr}
\toprule
Supervision & IoU $\uparrow$ & $L_2\downarrow$ & Overlap $\downarrow$ & \# prim. $\downarrow$\\
\midrule
SuperDec loss & 0.59 & 0.050 & 4\% & 5.78\\
$\mathcal L_{\rm IoU}$ only & \textbf{0.74} & 0.124 & 32\% & 5.94\\
$\mathcal L_{\rm IoU}+\mathcal L_{\rm SDF}$ & 0.72 & \textbf{0.042} & 14\% & 5.80\\
\rowcolor{pale}+ $\mathcal L_{\rm overlap}$ & 0.72 & 0.043 & \textbf{8\%} & \textbf{5.64}\\
\bottomrule
\end{tabular}\end{center}
\vspace{.14cm}
\begin{itemize}\setlength\itemsep{.42em}
\item \caccent{IoU alone} greatly improves volumetric reconstruction but produces excessive overlap.
\item \caccent{The SDF term} improves surface precision, reducing $L_2$ from 0.124 to 0.042.
\item \caccent{The overlap term} reduces redundancy from 14\% to 8\%.
\end{itemize}
\src{Paper Table 4. Sparsity and existence are held fixed; the IoU score does not improve monotonically with every added term.}
\end{frame}


% 19
\begin{frame}{Why keep deformation and test-time optimization?}
\begin{columns}[T,onlytextwidth]
\column{.49\textwidth}
\centering\textbf{Deformation ablation (refined)}\par\vspace{.15cm}
\begin{tabular}{lcc}
\toprule
Taper/Bend & IoU & $L_2$\\\midrule
Neither & 0.81 & 0.041\\
Taper only & 0.83 & 0.038\\
Bend only & 0.84 & 0.037\\
\rowcolor{pale}Both & \textbf{0.85} & \textbf{0.036}\\\bottomrule
\end{tabular}
\column{.49\textwidth}
\centering\textbf{Test-time optimization}\par\vspace{.15cm}
\begin{tabular}{lcc}
\toprule
 & No TTO & With TTO\\\midrule
IoU & 0.72 & \textbf{0.87}\\
F-score & 0.37 & \textbf{0.49}\\
Runtime & \textbf{0.0082 s} & 5 s\\\bottomrule
\end{tabular}
\end{columns}
\vspace{.23cm}\centering\caccent{Refinement improves geometry, learned feed-forward inference preserves speed.}
\end{frame}


% 20
\subsection{Partial point clouds}
\begin{frame}[shrink=3]{Occluded objects}
\begin{columns}[T,onlytextwidth]
\column{.51\textwidth}
\centering\safeincludegraphics[width=\linewidth,height=4.6cm,keepaspectratio]{assets/paper_scannet.png}
\src{Paper Fig. 7: qualitative ScanNet++ scene reconstructions.}
\column{.47\textwidth}
\vspace{.2cm}
\begin{tabular}{lr}
\toprule
ASE partial-cloud method & IoU\\\midrule
SuperDec & 0.20\\
SuperFlex, base & 0.20\\
+ occlusion fine-tuning & 0.48\\
\rowcolor{pale}+ ASE fine-tuning & \textbf{0.54}\\\bottomrule
\end{tabular}
\vspace{.23cm}
\begin{block}{Interpretation}
The model completes unseen geometry through priors learned from high-quality full-object decompositions. Better fitting on complete inputs alone does not solve occlusion.
\end{block}
\end{columns}
\end{frame}


% 21
\section{Conclusions}
\begin{frame}{Conclusions}
\vspace{.12cm}
{\large\bfseries SuperFlex brings geometric abstraction closer to accurate 3D reconstruction.}\par
\vspace{.23cm}
A complex object can be approximated by a small collection of explicit and interpretable parts. SuperFlex shows that this compact representation can preserve substantially more geometry than previous rigid superquadric decompositions.
\vspace{.22cm}
The main advance does not come from replacing the Transformer backbone. It comes from improving how primitives are fitted. The joint volumetric and surface loss produces more coherent reconstructions, while bending and tapering allow a single primitive to follow curved or asymmetric parts.
\vspace{.22cm}
Beyond complete point clouds, the method extends this representation to partially observed objects. The missing geometry is inferred from learned shape priors rather than measured directly.
\vspace{.22cm}
{\color{accent}\rule{\linewidth}{.55pt}}\par\vspace{.07cm}
\end{frame}
\appendix

\section{Appendix}
% Appendix 1
\begin{frame}{Appendix A -- Exact $z$-axis deformation equations}
\vspace{.1cm}
\begin{equation*}
T^{-1}(x,y,z)=\begin{pmatrix}\frac{x}{1+\tau_x z/s_z}\\[.2cm]\frac{y}{1+\tau_y z/s_z}\\ z\end{pmatrix}.
\end{equation*}
\vspace{.2cm}
\begin{equation*}
B_z^{-1}(x,y,z)=\begin{pmatrix}x-(A-r)\cos\alpha_z\\ y-(A-r)\sin\alpha_z\\\gamma/k_z\end{pmatrix},\quad
\begin{aligned}
A&=x\cos\alpha_z+y\sin\alpha_z,\\
\gamma&=\operatorname{atan2}(z,k_z^{-1}-A),\\
r&=k_z^{-1}-\sqrt{z^2+(k_z^{-1}-A)^2}.
\end{aligned}
\end{equation*}
\src{Paper Appendix A. Coordinate branches and $k=0$ limits need care; $\operatorname{atan2}$ makes the intended quadrant choice explicit.}
\end{frame}


% Appendix 2
\begin{frame}{Appendix B -- On the limits of radial distance}
\vspace{.13cm}
\begin{itemize}\setlength\itemsep{.4em}
\item A rigid superellipsoid satisfies $q(\lambda x)=\lambda^{2/\epsilon_1}q(x)$.
\item Exact ray intersection at $\lambda_{\ast}=q(x)^{-\epsilon_1/2}$; radial displacement $\|x\|(1-\lambda_{\ast})$.
\item Not the Euclidean SDF because nearest surface point need not lie on the center ray.
\item With nonlinear tapering/bending, $q(D^{-1}(\lambda x))$ generally is not homogeneous; substitution into the same closed form is a surrogate even for the radial ray intersection.
\end{itemize}
\begin{center}\safeincludegraphics[width=.67\linewidth,height=2.5cm,keepaspectratio]{assets/demo_radial.png}\end{center}
\src{Paper Eq. (3) describes this as radial distance. Geometric qualification follows from the lack of homogeneity after deformation.}
\end{frame}


% Appendix 4
\begin{frame}{Appendix C -- Matching and occlusion loss}
\vspace{.05cm}
\begin{equation*}
\mathcal L_{\rm geom}=\sum_{p\in\mathrm{matched\ active}}\operatorname{CD}\!\left(\mathcal S(\hat\Theta_p),\mathcal S(\Theta_p)\right).
\end{equation*}
\begin{equation*}
\boxed{\mathcal L_{\rm sup}=\mathcal L_{\rm geom}+\lambda_{\rm IoU}\mathcal L_{\rm IoU}+\mathcal L_{\rm rigid}}
\end{equation*}
\begin{itemize}\setlength\itemsep{.46em}
\item Use a Hungarian one-to-one assignment between the predicted and pseudo-ground-truth primitive sets.
\item $\mathcal L_{\rm rigid}$ compares counterparts with deformations switched off.
\item Necessary because different parameterizations, symmetries and permutations can produce similar surfaces.
\end{itemize}
\src{Paper Sec. 4.3 and Eq. (12).}
\end{frame}


% Appendix 6
\begin{frame}{Appendix D -- Implementation details}
\begin{columns}[T,onlytextwidth]
\column{.51\textwidth}
\begin{block}{Training choices}
\begin{tabular}{ll}
Point count $N$ & 4096\\
Primitive queries $P$ & 16\\
Feature width $D$ & 128\\
Decoder layers $L$ & 3\\
Epochs & 1000\\
Initial learning rate & $3\times10^{-4}$\\
Robust fine-tuning LR & $10^{-4}$\\
IoU temperature & $10^{-3}$\\
\end{tabular}
\end{block}
\column{.47\textwidth}
\begin{block}{Reproducing authors' inference}
\texttt{git clone https://github.com/}\newline\texttt{GabrielTavernini/superflex}\par\vspace{.17cm}
\texttt{pip install -r requirements.txt}\par\texttt{pip install -e .}\par\texttt{bash scripts/download\_checkpoints.sh}\par\texttt{python demo\_viser.py}
\end{block}
\end{columns}
\end{frame}


% Appendix 7
\begin{frame}{Appendix E -- From a single image to 3D primitives}
\begin{center}\safeincludegraphics[width=.83\linewidth,height=3.5cm,keepaspectratio]{assets/paper_rgb_pipeline.png}\end{center}
\begin{itemize}\setlength\itemsep{.35em}
\item SAM 3 masks instances in the original RGB image.
\item MoGE produces partial 3D point maps from a single view.
\item Robust SuperFlex turns partial object point sets into complete primitive abstractions.
\end{itemize}
\src{Paper supplementary Fig. 13. This is an application, not the core inference pipeline or an additional benchmark result.}
\end{frame}


% References
\begin{frame}{References and resources}
\begin{itemize}\setlength\itemsep{.65em}
\item Tavernini et al., \textit{SuperFlex: Deformable Superquadrics for Point Cloud Decomposition}, ECCV 2026, arXiv:2607.01015.
\item Fedele et al., \textit{SuperDec: 3D Scene Decomposition with Superquadric Primitives}, ICCV 2025.
\item Barr, \textit{Superquadrics and Angle-Preserving Transformations}, IEEE CG\&A, 1981.
\item Liu et al., \textit{Marching-Primitives: Shape Abstraction from Signed Distance Function}, CVPR 2023.
\item SuperFlex project: \texttt{https://superflex3d.github.io/}
\item Source and checkpoints: \texttt{https://github.com/GabrielTavernini/superflex}
\end{itemize}
\src{All source-paper figures included for scholarly presentation and credited on their slides. Notebook illustrations are original.}
\end{frame}


\end{document}